\section{Analysis: from contextual memory to first-contact control}\label{sec:theory}
The analysis follows the information available to the robot. First, the shared prior determines which unseen combinations can be inferred from previous experience. Second, a uniform memory-error bound quantifies what can be learned when context assignments remain uncertain. These two results supply the memory-accuracy term in the first-contact prediction bound. A simulation lemma then translates prediction accuracy into a decision gap. Recognition, probing, and safety results characterize how the same context belief is used after information becomes available.

\subsection{Assumptions and scope}
Let $x=(z,a)\in\mathcal Z_A$ denote a state--action input in the domain of analysis. Write $c_t$ for the true context at transition $t$, suppressing episode indices. The memory-confidence result assumes stationary linear realizability,
\begin{equation}\label{eq:realizability}
 \mu_c(x)=UW_c\phi(x),\qquad v_t=W_{c_t}\phi_t+\xi_t,
\end{equation}
with fixed $U$, $\phi$, and $W_c$. Conditional on the history immediately before $v_t$, each $\xi_{t,i}$ is zero-mean and $\sigma_{i,t}$-sub-Gaussian. This history includes the chosen state and action, available cues, and the true current context for analysis; the robot does not observe that context. Contexts and actions may depend on earlier observations but not on the current noise. The noise bound $\sigma_{i,t}=\sigma_i(z_t,a_t)>0$ and update weights are measurable before that noise is observed.

The Gaussian shared-prior model is an additional assumption for Proposition~\ref{prop:composition}, not a consequence of linear realizability. The frequentist confidence result instead retains explicit prior mismatch. Define the finite separation
\begin{equation}\label{eq:separation}
 \Delta_{\max}=\sup_{c,c',x\in\mathcal Z_A}
 \|\mu_c(x)-\mu_{c'}(x)\|_2.
\end{equation}
Realizability, stationary weights, noise domination, and bounded input domains delimit the confidence analysis. Misspecified features, underestimated noise, or unmodeled drift require additional error terms. The cost analysis uses the noise-free conditional mean map $F_c(x)=f(x)+\mu_c(x)$. The cost and safety results state their own regularity and feasibility assumptions.

\subsection{When unseen combinations are identifiable}
\label{sec:composition-theory}

The context kernel separates two questions: whether the observed combinations identify the additive contribution of a new combination, and whether that combination has an interaction that no other context reveals. The following result makes this distinction exact in the limit in which the relevant stored predictions are known without error.

\begin{proposition}[Identifiability of an unseen combination]
\label{prop:composition}
Fix one output coordinate and one feature vector $\phi$, and let $g(c)=\phi^\top w_c$ and $\chi=\phi^\top\Sigma_0\phi>0$. Assume $\kappa_0^2>0$ and $\kappa_j^2>0$ for every factor. Under the context prior,
\[
g(c)=\psi(c)^\top b+\xi_c,
\]
where $\psi(c)$ contains an intercept and the one-hot indicators of each factor value, $b$ has a positive definite diagonal Gaussian covariance, and the $\xi_c\sim\mathcal N(0,\kappa_\times^2\chi)$ are independent of $b$ and one another. Let $\Omega$ be a set of contexts with exactly observed $g(\Omega)$, including the training context, and let $q\notin\Omega$ use factor values represented in $\Omega$. With $K_\Omega=[k_{\mathcal C}(c,c')]_{c,c'\in\Omega}$,
\begin{align}
\operatorname{Var}\bigl(g(q)\mid g(\Omega)\bigr)&=\chi s_q,\label{eq:composition-variance}\\
s_q&=k_{\mathcal C}(q,q)-k_{\mathcal C}(q,\Omega)K_\Omega^\dagger k_{\mathcal C}(\Omega,q),\nonumber
\end{align}
where $\dagger$ denotes the Moore--Penrose inverse.
\begin{enumerate}[label=(\roman*)]
\item If $\kappa_\times^2=0$, the conditional variance is zero if and only if
\[
\psi(q)\in\operatorname{span}\{\psi(c):c\in\Omega\}.
\]
For two factors, this holds if and only if the two factor values of $q$ belong to the same connected component of the bipartite graph whose observed combinations are edges.
\item If $\kappa_\times^2>0$, then
\begin{align*}
\operatorname{Var}\bigl(g(q)\mid g(\Omega)\bigr)
&=\kappa_\times^2\chi
+\operatorname{Var}\bigl(\psi(q)^\top b\mid g(\Omega)\bigr)\\
&\geq\kappa_\times^2\chi.
\end{align*}
\end{enumerate}
\end{proposition}

Connectivity is the relevant requirement for two-factor transfer. Observing each factor value somewhere is insufficient: observations in disconnected components do not determine their relative additive contributions. A connected chain of combinations supplies those missing relations. More observations can resolve uncertainty about the additive terms, but cannot reveal an independent interaction at a combination that has never occurred.

Equation~\eqref{eq:composition-variance} is a conditional variance, not a realized prediction error. Under the specified Gaussian model it equals the conditional mean-squared error of the conditional-mean predictor at the chosen input. A particular error may be larger or smaller. Exact recovery additionally requires an additive model and exact stored observations; finite memory uncertainty leaves additional uncertainty. The pseudoinverse is necessary because the purely additive kernel can be singular even when all its additive variance parameters are positive. If $\chi=0$, the scalar prediction is deterministic under this prior and the span criterion is unnecessary. Fitted kernel parameters and approximate conditioning on uncertain memories provide operational estimates of this variance; they do not, by themselves, certify calibration of the fitted prior.

\subsection{Memory accuracy with uncertain context assignments}
\label{sec:confidence-theory}

Predictable assignment weights preserve the conditional zero-mean property of observation noise. They do not ensure that a sample comes from the memory being updated. The next result separates the resulting cross-context contamination from estimation noise and prior mismatch. It applies to an instantiated memory representing a stationary context; an unresolved new-value hypothesis that aggregates several contexts need not satisfy this condition.

For memory $c$, choose a positive definite row precision $P_{c,i}$ at creation and keep it fixed. It may depend on the information available at creation, but the sufficient statistics below contain only subsequent observations. In particular, the result does not cover selecting and replaying earlier residuals using their observed values. Write
\begin{align}
\omega_{c,i,s}&=\frac{q_s(c)}{\sigma_{i,s}^2},\qquad
V_{c,i,t}=P_{c,i}+\sum_{s<t}\omega_{c,i,s}\phi_s\phi_s^\top,\label{eq:weighted-design}\\
\widehat w_{c,i,t}&=V_{c,i,t}^{-1}\left(P_{c,i}\nu_{c,i,t}
+\sum_{s<t}\omega_{c,i,s}\phi_s v_{s,i}\right),\label{eq:weighted-estimate}\\
u_{c,i,t}(\phi)&=\sqrt{\phi^\top V_{c,i,t}^{-1}\phi}.\nonumber
\end{align}
Weights are zero before creation. The predictive probability $q_s(c)$ includes the available episode cue but excludes the current residual $v_s$. The prior centre $\nu_{c,i,t}$ may be recomputed from the observation history; this flexibility does not extend to changing $P_{c,i}$ while retaining these sufficient statistics.

\begin{theorem}[Uniform memory error under context uncertainty]
\label{thm:confidence}
Under the stationary realizability and conditional sub-Gaussian noise assumptions, fix a memory $c$ and a failure probability $\delta_c\in(0,1)$. Conditional on the information at its creation, with probability at least $1-\delta_c$, simultaneously for every subsequent $t$, every $\phi\in\mathbb R^p$, and every output coordinate $i\in\{1,\ldots,k\}$,
\begin{equation}
\left|\phi^\top(\widehat w_{c,i,t}-w_{c,i})\right|
\leq \beta_{c,i,t}\,u_{c,i,t}(\phi),\label{eq:memory-confidence}
\end{equation}
where
\begin{align}
\beta_{c,i,t}
&=\sqrt{2\log\frac{k}{\delta_c}
+\log\frac{\det V_{c,i,t}}{\det P_{c,i}}}
+\|\nu_{c,i,t}-w_{c,i}\|_{P_{c,i}}+B_{c,i,t},\label{eq:memory-radius}\\
B_{c,i,t}^2
&=\sum_{s<t}\omega_{c,i,s}
\left[\phi_s^\top(w_{c_s,i}-w_{c,i})\right]^2.\label{eq:contamination}
\end{align}
Here $c_s$ is the true context. Assigning failure probabilities with $\sum_c\delta_c\leq\delta$ makes these inequalities simultaneous over all instantiated memories with probability at least $1-\delta$.
\end{theorem}

If all observations lie in the domain defining $\Delta_{\max}$, the contamination is bounded by
\[
B_{c,i,t}^2\leq\Delta_{\max}^2 M_{c,i,t},\qquad
M_{c,i,t}=\sum_{s<t:c_s\ne c}\omega_{c,i,s}.
\]

The noise term follows the self-normalized martingale bound for predictable linear regression~\citep{abbasi2011improved}. Predictability is the essential requirement: the feature, noise scale, and assignment weight must be fixed before the current noise is observed. Neither independence between context assignments and earlier observations nor independent output coordinates is required. The simultaneous statement over inputs permits evaluation on the candidate trajectories considered by the planner. Fixed feature maps, fixed correction directions, and fixed per-memory precisions are part of the scope of this result.

The contamination term connects memory accuracy to context prediction. If every noise scale is at least $\sigma_{\min}>0$, then, for each output coordinate,
\begin{equation}
\sum_c M_{c,i,T}
\leq\frac{1}{\sigma_{\min}^2}\sum_{t<T}\bigl[1-q_t(c_t)\bigr]
\leq\frac{1}{\sigma_{\min}^2}\sum_{t<T}\ell_t^{\mathrm{dyn}},
\qquad \ell_t^{\mathrm{dyn}}=-\log q_t(c_t).
\label{eq:contamination-logloss}
\end{equation}
The sum on the left includes only instantiated memories. If the true context has no represented slot, set its probability to zero; the log-loss bound is then infinite and supplies no finite claim for that step. Unresolved new-value slots are not stationary memories covered by the theorem. Once contexts are represented, accurate predictions reduce both the error of selecting a memory and the contamination of memories updated with uncertain assignments.

Theorem~\ref{thm:confidence} is an error bound, not a claim that an uninflated Gaussian posterior interval has nominal coverage. Both the prior-mismatch norm and the contamination term involve unknown quantities. A computable certified radius requires valid upper bounds on these terms and on the noise scale; substituting fitted estimates does not preserve the coverage guarantee automatically. Known parameter ranges in state-based simulations can supply conservative envelopes. Empirical coverage evaluates fitted intervals when certified envelopes are unavailable. Negative transfer appears explicitly as a large prior-mismatch term. Additional data reduce $u_{c,i,t}(\phi)$ only in directions sufficiently explored by the weighted features, and persistent contamination need not vanish with sample size.

For use in the first-contact bound, on the simultaneous event a valid vector error envelope is
\[
\epsilon_t=\sup_{c,\,(z,a)\in\mathcal Z_A}
\left(\sum_{i=1}^k\beta_{c,i,t}^2
u_{c,i,t}\bigl(\phi(z,a)\bigr)^2\right)^{1/2},
\]
with the supremum restricted to the memories under consideration. Orthonormality of $U$ transfers this projected bound to the latent correction. A finite, useful envelope requires bounded prior mismatch, controlled contamination, and adequate feature coverage.

\paragraph{Why filtered weights can introduce bias.}
Consider scalar observations $v\sim\mathcal N(0,\sigma^2)$ and two fixed scoring models with means $0$ and $\Delta>0$, equal variances, and equal predictive probabilities. Hard assignment by the filtered probability sends an observation to the first model exactly when $v<\Delta/2$. A separate sample-mean estimator trained on those selected observations converges to
\[
\mathbb E[v\mid v<\Delta/2]
=-\sigma\frac{\varphi_{\mathrm N}(\Delta/(2\sigma))}
{\Phi_{\mathrm N}(\Delta/(2\sigma))}<0,
\]
where $\varphi_{\mathrm N}$ and $\Phi_{\mathrm N}$ are the standard normal density and distribution function. At $\Delta=2\sigma$, the bias is approximately $-0.29\sigma$. An interval that shrinks around this selected-sample mean eventually excludes zero. This counterexample keeps the scoring models fixed; it does not assert the same limiting value when those scoring models are learned online. Predictable weights remove this dependence on the current noise, while the contamination term accounts for samples generated by other contexts.


\subsection{First-contact error is controlled by memory accuracy and context log loss}
\label{sec:firstcontact-theory}
For a common query $x_e=(z_{e,\tau_e},a_{e,\tau_e})$, define the mean prediction bias
\[
 b_e^{\mathrm{pred}}=\left\|\sum_c\pi_e(c)m_{c,e}(x_e)-\mu_{c_e}(x_e)\right\|_2,
 \qquad \pi_e(c)=\pi_{e,\tau_e^-}(c).
\]
This quantity compares conditional means; it excludes irreducible observation noise. The error of an observed transition additionally includes that noise.

\begin{proposition}[First-contact separation]\label{prop:firstbound}
Consider a closed library of $K$ contexts containing each current context, and suppose
$\|m_{c,e}(x)-\mu_c(x)\|_2\leq\epsilon_e$ for all supported contexts and queries under consideration. Then
\begin{align}
 b_e^{\mathrm{pred}}&\leq\epsilon_e+
          \bigl(1-\pi_e(c_e)\bigr)\Delta_{\max},\label{eq:firstbound}\\
 \sum_{e=1}^N b_e^{\mathrm{pred}}
 &\leq\sum_{e=1}^N\epsilon_e+
       \Delta_{\max}\sum_{e=1}^N-\log\pi_e(c_e).\label{eq:firstlogloss}
\end{align}
By contrast, under Observation~\ref{obs:first}, an episode-reset predictor initialized at $f$ has bias $\|\mu_{c_e}(x_e)\|_2$. If this norm is at least $\Delta_{\min}>0$ on $N_{\mathrm{shift}}$ evaluated episodes, its cumulative bias is at least $N_{\mathrm{shift}}\Delta_{\min}$.
\end{proposition}
The upper bound for contextual prediction is smaller than the reset lower bound whenever
\begin{equation}\label{eq:separationcondition}
 \sum_e\epsilon_e+\Delta_{\max}\sum_e-\log\pi_e(c_e)
 <N_{\mathrm{shift}}\Delta_{\min}.
\end{equation}
This is a sufficient condition for separation on the specified queries, not a guarantee for every deployment. Theorem~\ref{thm:confidence} supplies $\epsilon_e$ on its simultaneous confidence event when the inputs and contexts satisfy its assumptions. The context log loss has two roles: episode-level loss controls anticipatory selection, while transition-level loss controls contamination in Eq.~\eqref{eq:contamination-logloss}.

The full method also retains a novel-context branch. If $p_\star$ is its aggregate mass, $m_\star$ its conditional mean, and the true context is already known, the corresponding one-query bound is
\begin{equation}\label{eq:novelbranchbound}
 b_e^{\mathrm{pred}}\leq(1-p_\star)\epsilon_e+
 (1-p_\star-\pi_e(c_e))\Delta_{\max}
 +p_\star\|m_\star(x_e)-\mu_{c_e}(x_e)\|_2.
\end{equation}
Without a bound for the last term, a closed-library guarantee does not apply to the entire mixture. Truly novel conditions require an acquisition phase or a valid transferred prior; no fixed one-episode acquisition cost is assumed.

\begin{corollary}[Finite-alphabet entropy reference]\label{cor:entropy}
Suppose the context alphabet of size $K$ and cue alphabet $\mathcal Y$ are fixed, and each episode's context identity is revealed after its prediction. For each cue value, use a categorical predictor with symmetric Dirichlet-$1/2$ counts:
\[
 \pi_e^{\mathrm{KT}}(c\mid y)=
 \frac{n_{<e}(c,y)+1/2}{n_{<e}(y)+K/2}.
\]
For any context--cue sequence, its cumulative log loss satisfies
\begin{equation}\label{eq:entropyreference}
 \sum_{e=1}^N-\log\pi_e^{\mathrm{KT}}(c_e\mid y_e)
 \leq N\widehat H(C\mid Y)+O\!\left(K|\mathcal Y|\log(N+1)\right).
\end{equation}
Substitution into Eq.~\eqref{eq:firstlogloss} gives the corresponding first-contact bias bound for this reference predictor.
\end{corollary}
This is the categorical coding bound of Krichevsky and Trofimov~\citep{krichevsky1981performance}, stated separately because its feedback assumptions differ from unlabeled context discovery. \method{} receives no context labels, and its inferred assignments do not automatically inherit Eq.~\eqref{eq:entropyreference}. Empirical conditional entropy therefore diagnoses available cue information, while measured context log loss tests the inference procedure. Neither quantity alone determines prediction error without memory accuracy and residual separation.

\subsection{Recognition depends on initial evidence and dynamics separation}
For a known true context $c$ and one fixed alternative $c'$, let
\[
 \ell_0=\log\frac{p(c\mid D_{<e})}{p(c'\mid D_{<e})}
       +\log\frac{p_\eta(y_e\mid c)}{p_\eta(y_e\mid c')}
\]
be the pre-contact log odds. A diagnostic observation contributes a log-likelihood ratio in favor of $c$.
\begin{proposition}[Two-context recognition time]\label{prop:recognition}
Assume fixed likelihoods and independent, identically distributed diagnostic observations under $c$, with integrable log-likelihood increments of mean $D>0$. Let $T$ be the first time the accumulated log odds reach a threshold $\ell_\star$, and assume $\E[T]<\infty$ and integrable threshold overshoot $O_T$. If $\ell_0\geq\ell_\star$, then $T=0$. Otherwise,
\begin{equation}\label{eq:recognition}
 \E[T]=\frac{\ell_\star-\ell_0+\E[O_T]}{D}
 \geq\frac{\ell_\star-\ell_0}{D}.
\end{equation}
For correctly specified likelihoods, $D$ is their Kullback--Leibler divergence under the diagnostic observation distribution.
\end{proposition}
The result isolates how a cue can replace post-contact evidence. Larger initial log odds reduce the threshold deficit, while better separated dynamics increase the evidence rate. Repeated exposure can improve both cue associations and memory likelihoods. Occurrence counts alone need not produce savings when competing contexts recur equally often. Adaptive actions, changing memories, multiple alternatives, and persistence-based recognition criteria require additional analysis; the two-context experiment controls these factors to test Eq.~\eqref{eq:recognition}.

\subsection{From model accuracy to the committed decision}
Prediction error matters through the action selected before contact. We connect the two for a fixed context over a noise-free horizon of $H$ transitions. A candidate $a=(a_0,\ldots,a_{H-1})$ is a committed action sequence, and its cost is
\begin{equation}\label{eq:horizoncost}
J_e(a)=\sum_{t=0}^{H-1}\ell(z_t^a,a_t)+g(z_H^a),
\end{equation}
where $g=0$ permits a problem without terminal cost. Assume the true map $F_{c_e}$ is $L$-Lipschitz in state, uniformly over candidate actions, and $\ell(\cdot,a_t)$ and $g$ are $L_\ell$-Lipschitz. These conditions hold on a domain containing both true and predicted trajectories. Set $C_H=\sum_{j=0}^{H-1}L^j$.

\begin{lemma}[Simulation bound]\label{lem:simulation}
Fix a candidate sequence $a$ and let $z_t$ and $\widehat z_t$ be the noise-free trajectories of $F_{c_e}$ and a model $\widehat F$, respectively, from the same initial state. Under the preceding Lipschitz assumptions,
\begin{equation}\label{eq:simulationbound}
|J_e(a)-\widehat J(a)|
\le L_\ell C_H\sum_{t=0}^{H-1}
\bigl\|F_{c_e}(\widehat z_t,a_t)-\widehat F(\widehat z_t,a_t)\bigr\|.
\end{equation}
\end{lemma}
The discrepancy is evaluated along the model's trajectory. That trajectory is available during planning; bounding its discrepancy additionally requires a valid error bound. The statement concerns fixed action sequences. Feedback policies require corresponding regularity assumptions on their closed-loop dynamics and costs.

\begin{corollary}[First-contact decision gap]\label{cor:decision}
Let $\mathcal A_e$ be a finite, common candidate set and write $\pi_e(c)=\pi_{e,\tau_e^-}(c)$ for a belief over a closed library containing the true context $c_e$. Let $\widehat J_{e,c}(a)$ roll out $f+m_c$, and set
\begin{equation}
\widehat J_e(a)=\sum_c\pi_e(c)\widehat J_{e,c}(a),
\qquad
\widehat J_e(\widehat a_e)\le\min_{a\in\mathcal A_e}\widehat J_e(a)+\eta_{\mathrm{opt},e}.
\end{equation}
Suppose $\|m_c(x)-\mu_c(x)\|\le\bar\epsilon_e$ for every supported context and every input along its candidate rollouts. Suppose also that $\|\mu_c(x)-\mu_{c_e}(x)\|\le\Delta_{\max}$ there. Then
\begin{equation}\label{eq:decisiongap}
J_e(\widehat a_e)-\min_{a\in\mathcal A_e}J_e(a)
\le 2L_\ell C_HH\left[\bar\epsilon_e+
\bigl(1-\pi_e(c_e)\bigr)\Delta_{\max}\right]+\eta_{\mathrm{opt},e}.
\end{equation}
With common constants across episodes,
\begin{align}\label{eq:cumulativedecisiongap}
\sum_{e=1}^N\left[J_e(\widehat a_e)-\min_{a\in\mathcal A_e}J_e(a)\right]
&\le 2L_\ell C_HH\left[\sum_{e=1}^N\bar\epsilon_e+
\Delta_{\max}\sum_{e=1}^N-\log\pi_e(c_e)\right]
+\sum_{e=1}^N\eta_{\mathrm{opt},e}.
\end{align}
\end{corollary}
The comparison is with the best candidate, and the uniform error requirement includes candidates the planner does not select. The bound separates memory accuracy, context selection, and optimization error. It can be loose over long horizons: when $L=1$, $C_H=H$. Quantitative use therefore requires validated dynamics and cost regularity; the statement does not extend automatically to discontinuous contact costs or stochastic rollouts.

\subsection{When probing changes the decision}
A probe is useful when its observations improve the ensuing decision enough to cover its cost. To isolate this information value, fix the state, a finite action set $\mathcal A$, and integrable costs $J_\theta(a)$, where $\theta=(c,W)$ includes context and memory parameters. For belief $b$ and the Bayesian posterior $b_o$ after outcome $o$ of probe $p$, define
\begin{equation}\label{eq:theoryvoi}
V(b)=\min_{a\in\mathcal A}\E_b[J_\theta(a)],
\qquad
\operatorname{VOI}(p)=V(b)-\E_o[V(b_o)].
\end{equation}

\begin{proposition}[Information value and action agreement]\label{prop:voi}
Under these fixed-state, fixed-action-set conditions, every probe has $\operatorname{VOI}(p)\ge0$. If one action $a^\dagger$ minimizes $J_\theta$ for $b$-almost every $\theta$, then $\operatorname{VOI}(p)=0$ for every probe. Nevertheless, $I(\theta;o\mid p)$ can be positive: for a discrete belief this occurs whenever two states of positive probability have different probe-outcome distributions.
\end{proposition}
Action disagreement is necessary, but not sufficient, for positive decision value. With nonnegative probe cost $\kappa(p)$, a rule requiring $\operatorname{VOI}(p)>\kappa(p)$ never selects a probe under common action agreement. Information gain alone can reward observations that leave the optimal decision unchanged.

The context-reweighting score used in planning approximates this criterion by holding memory parameters and the initial state fixed. A physical nudge can also change the state, refine a memory, or change which strikes pass the safety filter. Those effects alter the decision problem and are outside Proposition~\ref{prop:voi}. Likewise, local reduction in the variance of a linearized cost difference measures decision-relevant uncertainty, but is not an exact value-of-information calculation.

\subsection{Reusing calibrated latent safety margins}
Calibration uses the full latent prediction error, including noise, context mistakes, and residuals outside the correction subspace. Let $\mathcal C_t^\delta$ be the nonempty plausible set and $\bar z'_{c,t}$ the prediction computed before observing $z_{t+1}$. Its score is $s_{c,t}=\|z_{t+1}-\bar z'_{c,t}\|$. With an $L_h$-Lipschitz barrier function $h$ and $0<\alpha\le1$, the filter checks
\begin{equation}\label{eq:theorysafetyconstraint}
h(\bar z'_{c,t})-L_hb_{c,t}\ge(1-\alpha)h(z_t),
\qquad c\in\mathcal C_t^\delta.
\end{equation}
Each step is attributed to exactly one memory: a covered memory satisfying $s_{c,t}\le b_{c,t}$ if any exists, and otherwise the most probable plausible memory. Ties are immaterial. Only its margin is updated, without projection, by
\begin{equation}\label{eq:theoryaci}
b_{c,t+1}=b_{c,t}+\gamma\bigl(\mathbf1\{s_{c,t}>b_{c,t}\}-\delta_s\bigr).
\end{equation}

\begin{proposition}[Long-run latent barrier violations]\label{prop:safety}
Assume Eq.~\eqref{eq:theorysafetyconstraint} is enforced exactly whenever feasible; otherwise a fallback satisfies $h(z_{t+1})\ge(1-\alpha)h(z_t)$. Suppose all calibration scores lie in $[0,\bar S]$, each margin is initialized in $[0,\bar S]$, $\gamma>0$, and $0<\delta_s<1$. Let $K_T$ count all calibration streams initialized during the first $T$ steps, including resets. With the attribution and update above, for every sequence of contexts and cues,
\begin{equation}\label{eq:safetybound}
\sum_{t=0}^{T-1}\mathbf1\{h(z_{t+1})<(1-\alpha)h(z_t)\}
\le\delta_sT+K_T\frac{\bar S+2\gamma}{\gamma}.
\end{equation}
\end{proposition}
This pathwise calibration bound follows the deterministic adaptive conformal argument~\citep{gibbs2021adaptive}; it does not require a correct context posterior or a Gaussian noise model. Its assumptions still require bounded scores and a safe fallback. An asymptotic violation rate of at most $\delta_s$ requires $K_T/T\to0$. Repeatedly resetting margins consumes the same allowance as creating additional streams.

The counted event is failure of the one-step latent barrier inequality. It is not the number of time steps spent outside the safe set, and a per-episode rate requires one calibrated decision per episode or an episode-level score. Physical interpretation depends on how $h$ represents the physical constraint. Finally, a quadratic program obtained by linearizing the filter must verify the nonlinear inequality, or include a valid remainder bound, to meet the exact-enforcement assumption.

